\begin{mistake}{2026-06-16}{02}

\begin{problemcontent}
设 \(f(x)\) 和 \(g(x)\) 在 \([0,1]\) 上连续，在 \((0,1)\) 内可导，且 \(f(0)=f(1)=-1\)，\(\int_0^1 f(x)\mathrm{d}x > \frac{1}{2}\)。试证：至少存在一点 \(\xi \in (0,1)\)，使 \(f'(\xi) + g'(\xi)[f(\xi) - \xi] = 1\)。
\end{problemcontent}

\begin{solutioncontent}
要证式移项为 \(f'(\xi) - 1 + g'(\xi)[f(\xi) - \xi] = 0\)。
令 \(F(x) = f(x) - x\)，原式结构变为 \(F'(x) + g'(x)F(x) = 0\)。
运用微分方程解法，两边同乘积分因子 \(e^{g(x)}\)，构造辅助函数：
\[ H(x) = F(x)e^{g(x)} = [f(x) - x]e^{g(x)} \]
接下来寻找 \(F(x)\) 的零点。由题意：
\[ F(0) = f(0) - 0 = -1 < 0 \]
\[ F(1) = f(1) - 1 = -2 < 0 \]
利用积分条件判断区间内部正负情况：
\[ \int_0^1 F(x) \mathrm{d}x = \int_0^1 [f(x) - x] \mathrm{d}x = \int_0^1 f(x) \mathrm{d}x - \frac{1}{2} > 0 \]
由积分中值定理，至少存在一点 \(c \in (0,1)\) 使得 \(F(c) = \int_0^1 F(x) \mathrm{d}x > 0\)。
此时 \(F(0)<0, F(c)>0, F(1)<0\)，根据零点定理，\((0,c)\) 和 \((c,1)\) 内分别存在一点 \(x_1, x_2\) 使得 \(F(x_1) = F(x_2) = 0\)。
因此 \(H(x_1) = H(x_2) = 0\)。对 \(H(x)\) 在 \([x_1, x_2]\) 上应用罗尔定理，存在 \(\xi \in (x_1, x_2) \subset (0,1)\) 使得 \(H'(\xi) = 0\)，即：
\[ [f'(\xi) - 1]e^{g(\xi)} + [f(\xi) - \xi]g'(\xi)e^{g(\xi)} = 0 \]
约去恒正的 \(e^{g(\xi)}\) 并移项，即得 \(f'(\xi) + g'(\xi)[f(\xi) - \xi] = 1\)。
\end{solutioncontent}

\mistakeknowledge{
\begin{itemize}
  \item \textbf{核心公式/定理}：辅助函数构造（微分方程法）：遇 \(F'(x) + p(x)F(x) = 0\)，构造 \(H(x) = F(x)e^{\int p(x)\mathrm{d}x}\)；积分中值定理、零点定理、罗尔定理。
  \item \textbf{易错点}：面对 \(f(0), f(1)\) 均为负的情况不知如何找零点。关键在于利用 \(\int_0^1 F(x)\mathrm{d}x > 0\) 结合积分中值定理，推出区间内部必存在一点 \(c\) 使得 \(F(c) > 0\)，从而形成“两负夹一正”寻找两个零点。
  \item \textbf{关联知识}：微分方程分离变量法在构造辅助函数中的应用：将目标式视作 \(y' + p(x)y = 0\)，分离变量积分得 \(ye^{\int p(x)\mathrm{d}x} = C\)，常数项等号另一侧即为辅助函数整体。
\end{itemize}}

\end{mistake}
